\section{Design Principles}

\begin{proposition}[Cheapest controllable variable principle]
When several barrier terms are elevated, the first intervention should target
the cheapest controllable variable that materially changes transition odds,
rather than the most morally impressive variable.
\end{proposition}

\begin{discussion}
This principle follows the logic of control design. In a recurrent failure,
raising internal resolve may be possible, but it is usually expensive, noisy,
and hard to verify. By contrast, moving the phone away, pre-marking the next
problem, or attaching the task to an already recurring boundary changes the
state-transition conditions directly and cheaply. In dynamical-systems terms,
small local manipulations can matter if they change the system near a sensitive
entry point.
\end{discussion}

\begin{proposition}[History-dependent reinforcement principle]
An intervention should be judged not only by whether it succeeds once, but by
whether it makes the next attempt easier through accumulated local history.
\end{proposition}

\begin{discussion}
This principle is where emergence enters practice. If a reader repeatedly uses
the same location, same startup cue, and same minimum viable entry action, then
future transitions inherit a partially trained pathway. That is not yet a proof
of full attractor formation in the strict mathematical sense, but it is a
reasonable design inference from medium-lane ideas about metastability, causal
emergence, and path dependence: repeated local structure can produce more
stable higher-level patterns.
\end{discussion}
